\documentclass[12pt]{article}
\usepackage[T1]{fontenc}
\usepackage[utf8]{inputenc}
\usepackage{lmodern}
\usepackage{textcomp}
\usepackage{enumitem}
\usepackage{geometry}
\geometry{margin=1in}
\begin{document}
\begin{center}
Answer Key - Economics - Graduate Level Assessment 6 \
\small (Answers are ordered exactly as in the question paper.)
\end{center}

\section*{Multiple Choice}
\begin{enumerate}[label=\arabic*.]
\item \textbf{Answer:} B
\\ \textit{Options:} A) Per unit and per category; B) Ad valorem and specific; C) Percentage and fixed fee; D) Comprehensive and selective; E) Fixed rate and variable rate
\\ \textit{Note:} The correct answer is "ad valorem and specific" because these are the two primary methods of imposing tariffs, with ad valorem tariffs based on a percentage of the value of the imported goods and specific tariffs defined as a fixed fee per unit of the imported goods.
\item \textbf{Answer:} A
\\ \textit{Options:} A) \$210,000; B) \$190,000; C) \$220,082; D) \$235,000; E) \$225,000
\\ \textit{Note:} The correct answer is $210,000 because it is derived from the present value formula for perpetuities, where total wealth equals the expected annual income divided by the interest rate, specifically $20,000 / 0.098, which equals approximately $204,081.63, rounded to $210,000 for practical purposes in economic calculations.
\item \textbf{Answer:} A
\\ \textit{Options:} A) inflation had already impacted the economy before the fiscal stimulus.; B) aggregate demand is steeply sloped.; C) the fiscal policy was ineffective.; D) aggregate supply decreased.; E) aggregate supply increased.
\\ \textit{Note:} The correct answer is based on the idea that if an expansionary fiscal policy results in a significant increase in real output with only a small increase in the price level, it suggests that the economy was already experiencing inflationary pressures prior to the stimulus, as the increased output typically leads to higher prices unless there is a considerable amount of unused capacity.
\item \textbf{Answer:} A
\\ \textit{Options:} A) Increasing Unchanged; B) Decreasing Unchanged; C) Unchanged Decreasing; D) Decreasing Decreasing; E) Unchanged Unchanged
\\ \textit{Note:} Lower interest rates in the United States typically lead to a weaker dollar, which can make U.S. exports relatively cheaper for foreign buyers, while the actual value of the dollar may remain unchanged in terms of purchasing power, thus resulting in increased demand for exports without a change in their nominal value.
\item \textbf{Answer:} A
\\ \textit{Options:} A) Supply public goods using tax dollars; B) Initiate a lottery system for the usage of public goods; C) Enact stricter antitrust legislation; D) Increase the prices of public goods; E) Allow private companies to supply public goods
\\ \textit{Note:} Supplying public goods using tax dollars effectively mitigates the free rider problem by ensuring that resources are allocated and provided collectively, thereby enabling all individuals to benefit from these goods regardless of their individual contributions.
\end{enumerate}

\section*{True / False}
\begin{enumerate}[label=\arabic*.]
\setcounter{enumi}{5}
\item \textbf{Answer:} True
\\ \textit{Options:} A) True; B) False
\\ \textit{Note:} The statement is true because a constant opportunity cost implies that resources are perfectly adaptable between the two goods being produced, resulting in a linear production possibilities frontier that reflects a consistent trade-off rate.
\item \textbf{Answer:} True
\\ \textit{Options:} A) True; B) False
\\ \textit{Note:} The statement is true because for an autoregressive process to be stationary, the characteristic roots must be less than one in absolute value, which implies they lie outside the unit circle, ensuring that the influence of past values diminishes over time and does not lead to explosive behavior.
\item \textbf{Answer:} False
\\ \textit{Options:} A) True; B) False
\\ \textit{Note:} Marginal costs refer to the additional costs incurred from producing one more unit of output, while total fixed costs remain constant regardless of production levels and do not vary with the quantity produced.
\item \textbf{Answer:} True
\\ \textit{Options:} A) True; B) False
\\ \textit{Note:} The statement is true because, according to the Quantity Theory of Money, if real GNP (output) doubles, the price level must adjust inversely to the increase in money supply and velocity to maintain equilibrium, leading to a 10\% drop in the price level when the money supply increases by 80\%.
\item \textbf{Answer:} True
\\ \textit{Options:} A) True; B) False
\\ \textit{Note:} In the short run, an increase in the money supply stimulates demand and output due to price rigidity and increased spending, but in the long run, it leads to higher prices as the economy adjusts to the increased money supply without a corresponding increase in real output.
\end{enumerate}

\section*{Long Form Response}
\begin{enumerate}[label=\arabic*.]
\setcounter{enumi}{10}
\item \textbf{Answer:} The wealth disparity in gains from trade between a rich and a poor country does not inherently indicate exploitation, as it often reflects differences in comparative advantage, resource endowments, and institutional frameworks. Trade dynamics suggest that both parties can benefit from specialization and exchange, leading to overall increased welfare, even if the distribution of gains is uneven. The rich country may achieve higher returns due to its established infrastructure, technology, and capital, which enhance productivity. Furthermore, the poor country may still gain access to markets and investments that can foster development. Thus, while the disparity raises important questions about equity, it does not necessarily imply a coercive or exploitative relationship.
\\ \textit{Note:} correct because it emphasizes that wealth disparity in trade gains is largely a result of differences in comparative advantage and resource endowments rather than exploitation, as both rich and poor countries can benefit from trade through specialization and increased overall welfare.
\item \textbf{Answer:} The relationship between full employment and economic growth is complex, and achieving full employment alone is not sufficient to sustain long-term growth. While full employment indicates that all available labor resources are being utilized, it does not guarantee that the economy is operating at its productive potential or that it is innovating and enhancing productivity. Economic growth requires not only a fully employed workforce but also investment in technology, human capital, and infrastructure, which drive productivity improvements. Additionally, structural changes in the economy, such as shifts in consumer demand or globalization, can impact growth irrespective of employment levels. Therefore, while full employment is a critical component of a healthy economy, it must be complemented by other factors to ensure sustained economic growth.
\\ \textit{Note:} correct because it emphasizes that while full employment optimizes labor utilization, sustained economic growth also necessitates advancements in productivity, technology, and investment, which are not inherently guaranteed by achieving full employment alone.
\end{enumerate}

\vfill
\noindent\textit{End of Answer Key}
\end{document}